\documentclass[10pt]{book}
\usepackage[sectionbib]{natbib}
\usepackage{array,epsfig,fancyhdr,rotating}
\usepackage{hyperref}
\usepackage{xr-hyper}
\externaldocument{profile_sieve_geometric_paper}
\usepackage{amsmath}
\usepackage{amssymb}
\usepackage{amsfonts}
\usepackage{amsthm}
\usepackage{bm}
\usepackage{graphicx}
\graphicspath{{writing_samples/figures/}{figures/}}

\textwidth=31.9pc
\textheight=46.5pc
\oddsidemargin=1pc
\evensidemargin=1pc
\headsep=15pt
\headheight=28pt
\topmargin=.6cm
\parindent=1.7pc
\parskip=0pt

\numberwithin{equation}{section}
\setcounter{page}{1}
\newtheorem{theorem}{Theorem}
\newtheorem{lemma}{Lemma}
\newtheorem{corollary}{Corollary}
\newtheorem{proposition}{Proposition}
\newtheorem{assumption}{Assumption}
\theoremstyle{definition}
\newtheorem{definition}{Definition}
\newtheorem{example}{Example}
\newtheorem{remark}{Remark}

\newcommand{\ind}{\mathbf 1}
\newcommand{\R}{\mathbb R}

\pagestyle{fancy}
\def\n{\noindent}
\lhead[\fancyplain{} \leftmark]{}
\chead[]{}
\rhead[]{\fancyplain{}\rightmark}
\cfoot{\thepage}

\begin{document}
\raggedbottom

\renewcommand{\baselinestretch}{2}

\markright{ \hbox{\footnotesize\rm Statistica Sinica: Supplement
%{\footnotesize\bf 24} (201?), 000-000
}\hfill\\[-13pt]
\hbox{\footnotesize\rm
}\hfill }

\markboth{\hfill{\footnotesize\rm SIYAN DONG} \hfill}
{\hfill {\footnotesize\rm MULTIPLE NUISANCE BREAKS AND PREDICTABILITY} \hfill}

\renewcommand{\thefootnote}{}
$\ $\par \fontsize{12}{14pt plus.8pt minus .6pt}\selectfont

\centerline{\large\bf MULTIPLE NUISANCE BREAKS AND}
\vspace{2pt}
\centerline{\large\bf SPURIOUS PREDICTABILITY UNDER}
\vspace{2pt}
\centerline{\large\bf PERSISTENT REGRESSORS}
\vspace{.25cm}
\centerline{Siyan Dong}
\vspace{.4cm}
\centerline{\it KU}
\vspace{.55cm}
\centerline{\bf Supplementary Material}
\vspace{.55cm}
\fontsize{9}{11.5pt plus.8pt minus .6pt}\selectfont
\noindent
This supplement gives the technical details behind the block-sieve empirical likelihood construction in the main manuscript. Section S1 records the finite-sample block projection identities. Section S2 proves the stable-score decomposition and the omitted-break failure theorem. Section S3 proves the known-partition Wilks theorem by extending the original profile-sieve geometry to a regime-specific nuisance space. Section S4 gives the estimated-partition equivalence argument based on the set of misclassified observations. Section S5 proves the local-alternative calculation. Section S6 records implementation details for the reported simulations and the empirical application.
\par

\setcounter{section}{0}
\setcounter{equation}{0}
\renewcommand{\thesection}{S\arabic{section}}
\renewcommand{\theequation}{\thesection.\arabic{equation}}
\renewcommand{\thelemma}{S\arabic{lemma}}
\renewcommand{\theproposition}{S\arabic{proposition}}
\renewcommand{\thecorollary}{S\arabic{corollary}}
\renewcommand{\thetheorem}{S\arabic{theorem}}
\renewcommand{\theremark}{S\arabic{remark}}

\fontsize{12}{14pt plus.8pt minus .6pt}\selectfont

\section{Block Projection Algebra}

Fix a partition \(\Lambda=(k_1,\ldots,k_q)\). Let \(\bm Q_{K,\Lambda}\) and \(\bm M_{K,\Lambda}\) be defined as in (\ref*{eq:block-q}) and (\ref*{eq:block-m}). Throughout this supplement, all projection equalities are equalities in the empirical Hilbert space \((\R^T,T^{-1}\bm a'\bm b)\).

\begin{lemma}[Block projection identities]
\label{lem:s-block-identities}
If \(\bm Q_{K,\Lambda}'\bm Q_{K,\Lambda}\) is nonsingular, then
\[
    \bm Q_{K,\Lambda}'\bm M_{K,\Lambda}=\bm0,
    \qquad
    \bm M_{K,\Lambda}'=\bm M_{K,\Lambda},
    \qquad
    \bm M_{K,\Lambda}^2=\bm M_{K,\Lambda}.
\]
For every \(\bm c\) conformable with \(\bm Q_{K,\Lambda}\),
\[
    (\bm Q_{K,\Lambda}\bm c)'\bm M_{K,\Lambda}\bm w=0.
\]
\end{lemma}

\begin{proof}
Let \(\bm\Pi_{K,\Lambda}=\bm Q_{K,\Lambda}(\bm Q_{K,\Lambda}'\bm Q_{K,\Lambda})^{-1}\bm Q_{K,\Lambda}'\). This is the orthogonal projector onto the column space of \(\bm Q_{K,\Lambda}\). Therefore \(\bm\Pi_{K,\Lambda}'=\bm\Pi_{K,\Lambda}\), \(\bm\Pi_{K,\Lambda}^2=\bm\Pi_{K,\Lambda}\), and \(\bm Q_{K,\Lambda}'(\bm I_T-\bm\Pi_{K,\Lambda})=\bm0\). Since \(\bm M_{K,\Lambda}=\bm I_T-\bm\Pi_{K,\Lambda}\), the displayed identities follow. The final equality is obtained by multiplying \(\bm Q_{K,\Lambda}'\bm M_{K,\Lambda}=\bm0\) by \(\bm c\) and \(\bm w\).
\end{proof}

\section{Stable-Score Failure}

This section proves Proposition \ref*{prop:stable-decomp} and Theorem \ref*{thm:failure}. Let \(\Lambda=0\) denote the no-break partition and write \(\bm M_{K,0}\) for the corresponding stable-nuisance residual-maker.

\begin{proof}[Proof of Proposition \ref*{prop:stable-decomp}]
Under \(H_0:\beta=\beta_0\),
\[
    \bm Y-\beta_0\bm X_{lag}=\bm\mu+\bm U.
\]
Multiplying by the stable residualized weight gives the exact identity
\[
    T^{-1/2}(\bm Y-\beta_0\bm X_{lag})'\bm M_{K,0}\bm w
    =
    T^{-1/2}\bm U'\bm M_{K,0}\bm w
    +
    T^{-1/2}\bm\mu'\bm M_{K,0}\bm w.
\]
The second term is the omitted-break component. No stable-sieve approximation remainder is needed because \(B_T\) is defined using the exact nuisance vector \(\bm\mu\).
\end{proof}

\begin{proof}[Proof of Theorem \ref*{thm:failure}]
The scalar EL expansion gives
\[
    \ell_T^{\mathrm{st}}(\beta_0)
    =
    \frac{\{S_T^{\mathrm{st}}(\beta_0)\}^2}{V_T^{\mathrm{st}}}
    +o_p(1).
\]
By Proposition \ref*{prop:stable-decomp}, \(S_T^{\mathrm{st}}(\beta_0)=A_T^{\mathrm{st}}+B_T\). If \(|B_T|/(V_T^{\mathrm{st}})^{1/2}\to_p\infty\), the quadratic ratio diverges in probability, so rejection at any fixed critical value has probability tending to one. If the normalized oracle term converges to \(N(0,1)\) and the normalized break term converges to a nondegenerate random variable \(B\), then the limiting quadratic form is \((Z+B)^2\), where \(Z\sim N(0,1)\) conditionally on the limiting variance. This is central chi-square only when \(B=0\) with probability one.
\end{proof}

\section{Known-Partition Wilks Theorem}

Set \(\Lambda=\Lambda_0\) and write
\[
    \bm v_0=\bm M_{K,\Lambda_0}\bm w,
    \qquad
    \widehat{\bm u}_0(\beta_0)
    =\bm M_{K,\Lambda_0}(\bm Y-\beta_0\bm X_{lag}).
\]
Under the true partition,
\[
    \bm Y-\beta_0\bm X_{lag}
    =
    \bm Q_{K,\Lambda_0}\bm c_0+\bm r_0+\bm U,
\]
where \(\bm r_0\) stacks the regime-specific sieve approximation errors.

\begin{lemma}[Known-partition replacement]
\label{lem:s-known-replacement}
Under Assumptions \ref*{ass:block-sieve} and \ref*{ass:oracle},
\begin{align}
    T^{-1/2}\sum_{t=1}^T\widehat Z_t(\beta_0,\Lambda_0)
    &=T^{-1/2}\sum_{t=1}^T U_tv_{t,0}+o_p(1),
    \label{eq:s-num}\\
    T^{-1}\sum_{t=1}^T\widehat Z_t(\beta_0,\Lambda_0)^2
    &=T^{-1}\sum_{t=1}^T U_t^2v_{t,0}^2+o_p(1),
    \label{eq:s-den}\\
    \max_{t\leq T}|\widehat Z_t(\beta_0,\Lambda_0)|&=o_p(\sqrt T).
    \label{eq:s-max}
\end{align}
\end{lemma}

\begin{proof}
Using Lemma \ref{lem:s-block-identities},
\[
    \widehat{\bm u}_0(\beta_0)
    =
    \bm M_{K,\Lambda_0}(\bm U+\bm r_0),
    \qquad
    \bm v_0=\bm M_{K,\Lambda_0}\bm w.
\]
Because \(\bm M_{K,\Lambda_0}\) is symmetric and idempotent,
\[
    T^{-1/2}\widehat{\bm u}_0(\beta_0)'\bm v_0
    =
    T^{-1/2}(\bm U+\bm r_0)'\bm v_0.
\]
The approximation term is bounded by
\[
    |T^{-1/2}\bm r_0'\bm v_0|
    \leq
    \sqrt T\|\bm r_0\|_T\|\bm v_0\|_T=o_p(1),
\]
which proves (\ref{eq:s-num}). For the denominator, write \(\widetilde{\bm U}_0=\bm M_{K,\Lambda_0}\bm U\). The feasible score decomposes as
\[
    \widehat Z_t(\beta_0,\Lambda_0)
    =
    \widetilde U_{0,t}v_{t,0}
    +
    (\bm M_{K,\Lambda_0}\bm r_0)_tv_{t,0}.
\]
The approximation contribution is negligible by Assumption \ref*{ass:block-sieve}. Assumption \ref*{ass:oracle} then gives
\[
    T^{-1}\sum_{t=1}^T
    \{(\widetilde U_{0,t}v_{t,0})^2-(U_tv_{t,0})^2\}=o_p(1),
\]
which proves the denominator equivalence (\ref{eq:s-den}). The maximum bound follows from the oracle maximum condition, the projected-noise replacement condition, and the sup-norm sieve approximation bound.
\end{proof}

\begin{lemma}[EL expansion]
\label{lem:s-el-expansion}
Let \(Z_{t,T}\) be scalar scores. Suppose the convex-hull condition holds and
\[
    T^{-1/2}\sum_t Z_{t,T}=O_p(1),
    \qquad
    T^{-1}\sum_t Z_{t,T}^2\to_p\eta^2,
    \qquad
    \max_t|Z_{t,T}|=o_p(\sqrt T),
\]
where \(P(\eta^2>\underline\eta)>1-o(1)\). Then
\[
    2\sum_{t=1}^T\log(1+\widehat\lambda Z_{t,T})
    =
    \frac{\{T^{-1/2}\sum_tZ_{t,T}\}^2}
         {T^{-1}\sum_tZ_{t,T}^2}
    +o_p(1).
\]
\end{lemma}

\begin{proof}
The multiplier equation implies \(\widehat\lambda=O_p(T^{-1/2})\). Since \(\max_t|\widehat\lambda Z_{t,T}|=o_p(1)\), a Taylor expansion of the multiplier equation gives
\[
    \widehat\lambda
    =
    \frac{\sum_t Z_{t,T}}{\sum_t Z_{t,T}^2}
    +o_p(T^{-1/2}).
\]
Substitution into the second-order expansion of the log empirical likelihood gives the displayed quadratic approximation.
\end{proof}

\begin{proof}[Proof of Theorem \ref*{thm:known-wilks}]
Apply Lemma \ref{lem:s-known-replacement} to reduce the feasible score to the matched block-projection oracle score \(U_tv_{t,0}\). Assumption \ref*{ass:oracle} gives the mixed-normal numerator limit, denominator consistency, and maximum-score condition. Lemma \ref{lem:s-el-expansion} then yields
\[
    \ell_T(\beta_0,\Lambda_0)
    =
    \frac{\{T^{-1/2}\sum_tU_tv_{t,0}\}^2}
         {T^{-1}\sum_tU_t^2v_{t,0}^2}
    +o_p(1)
    \Rightarrow\chi_1^2.
\]
\end{proof}

\section{Estimated-Partition Equivalence}

This section records the high-level argument behind Theorem \ref*{thm:estimated-equiv}. Let \(\widehat\Lambda=\widehat\Lambda(\beta_0)\) and define
\[
    \mathcal M_T=
    \{t:j_{\widehat\Lambda}(t)\neq j_{\Lambda_0}(t)\}.
\]
Assumption \ref*{ass:partition} states that \(|\mathcal M_T|=o_p(\sqrt T)\) and that the boundary observations are negligible at the numerator and denominator scales.

\begin{lemma}[Misclassification bound]
\label{lem:s-misclassified}
Suppose Assumption \ref*{ass:partition} holds. Suppose also that the projected weights are uniformly bounded in probability, the score arrays satisfy the same maximum-score bounds as in Assumption \ref*{ass:oracle}, and the segment Gram matrices are stable uniformly over admissible partitions. Then
\begin{align*}
    T^{-1/2}\sum_{t=1}^T
    \{\widehat Z_t(\beta_0,\widehat\Lambda)-\widehat Z_t(\beta_0,\Lambda_0)\}
    &=o_p(1),\\
    T^{-1}\sum_{t=1}^T
    \{\widehat Z_t(\beta_0,\widehat\Lambda)^2-\widehat Z_t(\beta_0,\Lambda_0)^2\}
    &=o_p(1),\\
    \max_t|\widehat Z_t(\beta_0,\widehat\Lambda)|&=o_p(\sqrt T).
\end{align*}
\end{lemma}

\begin{proof}
Away from \(\mathcal M_T\), the estimated partition and true partition assign observations to the same regime. The two residual makers differ only through least-squares fits perturbed by the boundary observations and by terms controlled by uniform Gram stability. With \(\mathcal Z_T\) defined in Assumption \ref*{ass:partition}, the direct contribution of the boundary set to the numerator is bounded by
\[
    T^{-1/2}|\mathcal M_T|\mathcal Z_T=o_p(1).
\]
The denominator contribution is bounded by \(T^{-1}|\mathcal M_T|\mathcal Z_T^2=o_p(1)\). The remaining terms are the uniform Gram-stability perturbations of the least-squares coefficients and are of the same smaller order. Uniform Gram stability transfers the maximum-score condition from the true partition to the estimated partition.
\end{proof}

\begin{proof}[Proof of Theorem \ref*{thm:estimated-equiv}]
Lemma \ref{lem:s-misclassified} shows that the estimated-partition score and the true-partition score have the same first-order numerator, denominator, and maximum-score behavior. Applying Lemma \ref{lem:s-el-expansion} to both arrays yields
\[
    \ell_T\{\beta_0,\widehat\Lambda(\beta_0)\}
    -
    \ell_T(\beta_0,\Lambda_0)
    =o_p(1).
\]
The Wilks limit follows from Theorem \ref*{thm:known-wilks}. If \(P\{\widehat q(\beta_0)=q_0\}\to1\), the same argument holds on the event of correct order selection, proving Corollary \ref*{cor:qhat}.
\end{proof}

\section{Local Predictive Alternatives}

This section proves Theorem \ref*{thm:local-power}. Under the local alternative \(\beta_T=\beta_0+d_T\), the null-imposed residual satisfies
\[
    \bm Y-\beta_0\bm X_{lag}
    =
    \bm\mu+\bm U+d_T\bm X_{lag}.
\]
At the true partition, the block projection removes the regime-specific sieve approximation to \(\bm\mu\), so the score numerator has the expansion
\[
    T^{-1/2}\sum_{t=1}^T \widehat Z_t(\beta_0,\Lambda_0)
    =
    T^{-1/2}\sum_{t=1}^T U_tv_{t,0}
    +
    d_TT^{-1/2}\bm X_{lag}'\bm M_{K,\Lambda_0}\bm w
    +o_p(1).
\]
By the local drift condition in Theorem \ref*{thm:local-power}, the second term converges to \(\eta_q\delta_h\). The denominator remains \(\eta_q^2+o_p(1)\) because the local drift is at the testing scale and the nuisance-break component is removed by \(\bm M_{K,\Lambda_0}\). Hence the quadratic EL expansion gives a \(\chi_1^2(\delta_h^2)\) limit for the known partition.

Lemma \ref{lem:s-misclassified} transfers the same numerator, denominator, and maximum-score behavior to \(\widehat\Lambda(\beta_0)\), proving the estimated-partition version.
\section{Implementation Details}

For simulations, each replication uses the same null-imposed profile path as the theory:
\[
    \beta
    \to
    \widehat q(\beta),\widehat\Lambda(\beta)
    \to
    \bm M_{K,\widehat\Lambda(\beta)}
    \to
    \widehat Z_t(\beta)
    \to
    \ell_T^{BA}(\beta).
\]
The proof object is global profile-sieve segmentation. Computation may use multiscale candidate generation, followed by final selection with the profile-sieve RSS and the information criterion in (\ref*{eq:qhat}). The implementation must report the candidate-generation method separately from the statistical estimator, because SBS, NOT, and related devices are computational accelerators rather than the main contribution.

For empirical work, the stable and break-adaptive statistics are reported together. The stable statistic documents conventional evidence of predictability; the adaptive statistic documents whether that evidence survives multiple nuisance breaks. The empirical section also reports \(\widehat q\), break dates, regime-specific nuisance fits, an estimated break-induced score component, and residualized predictor strength \(T^{-1}\bm w'\bm M_{K,\widehat\Lambda}\bm w\).

\end{document}
